\paragraph{\trCountingSections。}\label{sub:counting-sections}
把\trFloer{} TQFT 的标准框架加以推广的一种方式，是把目标空间换成一个\trLefschetzFibration，并考察\trPseudoholomorphic\trSection{}而非映射。我们将从最简单的可能情形，即基\trSurface\trCompact{}的情形出发来解释这一推广，然后再扩大该框架以允许\trStriplikeEnd。我们这样做的意图，是要把代数的扭转构造 $T$ 与几何的\trDehnTwist{} $\tau$ 联系起来（\cref{thm:twist-equals-tau}）。我们会格外细心地以一种相当\trCanonical{}的形式来建立这一等同，使其随后可以轻松地过渡到 $\ZZ/2$-\trEquivariant{}的语境。
设 $\pi\colon E^{2n+2}\to S$ 是一个（照常\trExact{}的）\trLefschetzFibration。一个\emph{\trLagrange\trBoundaryCondition}（Lagrangian boundary condition）是一个\trSubmanifold{}$F^{n+1}\subset E$，它含于 $\partial^vE$ 之中且与 $\partial^hE$ 不相交，并满足：存在 $\alpha_F\in\Omega^1(\partial S)$ 与 $h_F\in C^\infty(F,\RR)$ 使得
\begin{equation}\label{eq:lagrangian-boundary-condition}
\theta_E|F=(\pi|F)^*(\alpha_F+\theta_S|\partial S)+dh_F.
\end{equation}
让我们把这稍微解读一下。$F$ 显然在 $E$ 中是\trLagrange{}的。由于 $\pi$ 在 $\partial^vE$ 中没有\trCriticalPoint，故 $\pi|F\colon F\to\partial S$ 必是一个\trFiberBundle（\cref{thm:critical-point-restriction}）。每个 $F_z\subset E_z$，$z\in\partial S$，实际上是\trExact{}的\trLagrange\trSubmanifold，而且这些\trSubmanifold{}会被沿边界的\trParallelTransport{}相互带入彼此（\cref{thm:fibered-lagrangian-criterion}）。在相当大的程度上，其逆亦真：设 $F_z\subset E_z$，$z\in\partial S$，是一族\trExact\trLagrange\trSubmanifold，它们被边界\trParallelTransport{}相互带入彼此，并令 $F$ 为它们的并。那么 $\theta_E|F$ 是\trClosed{}的，而且可以找到一个函数 $h_F$ 使得 $\theta_E|F-dh_F$ 限制在每个 $F_z$ 上都为零。若再假定诸 $F_z$ 都\trConnected，则 $\theta_E|F-dh_F$ 必然是 $\partial S$ 上某个 1-形式的\trPullback，从而 $F$ 是一个\trLagrange\trBoundaryCondition。
取定这样一个 $F$。一个\emph{\trRelative\trPerturbationDatum}（relative perturbation datum）是一对 $(K,J)$，其中：$K$ 是 $E$ 上的一个 1-形式，它在属于 $TE^v$ 的向量上为零，在 $\partial^hE\cup\mathrm{Crit}(\pi)$ 的某个邻域内整体为零，并且满足 $K|F=0\in\Omega^1(F)$；而 $J$ 是 $E$ 上的一个\trAlmostComplex{}结构，它$\phi_E$-\trCompatible、使 $\pi$ 成为\trPseudoholomorphic{}的，并且在 $\partial^hE\cup\mathrm{Crit}(\pi)$ 的某个邻域内与 $I_E$ 一致。在此前提下，一个边界落在 $F$ 中的\emph{\trInhomogeneous\trPseudoholomorphic\trSection}（inhomogeneous pseudo-holomorphic section）是一个满足如下条件的映射 $u\colon S\to E$：
\begin{equation}\label{eq:section-equation}
\begin{cases}
\pi(u(z))=z,\\
Du(z)+J(u)\circ Du(z)\circ I_S=Y(u)+J(u)\circ Y(u)\circ I_S,\\
u(\partial S)\subset F.
\end{cases}
\end{equation}
这里 $Y$ 是丛 $\Hom(\pi^*TS,TE^v)$ 的一个\trSection，由 $K$ 按下述方式确定：任给 $\xi\in TS_z$，$z=\pi(x)$，可以把它\trLift{}为 $E_z$ \trRegular{}部分上的一个切向量 $X$；函数 $K(X)$ 与\trLift{}的选取无关，而人们把 $Y(\xi)_x$ 定义为 $E_z$ 上相应的\trHamiltonianVectorField{}在 $x$ 处的取值。由于要求 $K$ 在\trCriticalPoint{}附近为零，$Y$ 亦然，从而 $Y$ 可以在这些点上平凡地延拓。
$Y$ 在\trCriticalPoint{}附近为零（从而在这些点上连续延拓），在 $\partial^hE$ 附近亦然。
\begin{remark}\label{thm:relative-perturbation-geometry}
\trRelative\trPerturbationDatum{}有一个简单的几何解释。考察修改后的（仍\trClosed{}且逐\trFiber{}\trSymplectic{}的）2-形式
\begin{equation}\label{eq:modified-form}
\phi_E^{\mathrm{new}}=\phi_E-dK
\end{equation}
及其相伴的\trConnection，它为
\begin{equation}\label{eq:modified-connection}
TE_x^{h,\mathrm{new}}=\{X=X^h+X^v\in TE_x=TE_x^h\oplus TE_x^v : X^v=Y(D\pi(X))\}.
\end{equation}
关于\trAlmostComplex{}结构，注意由假设 $J$ 保持\trSplitting{}$TE^h\oplus TE^v$，并且在第一个加项上等于 $I_E$。因此，这里的选择本质上就是 $TE^v$ 上的一个\trAlmostComplex{}结构{}的选择。存在唯一的 $J^{\mathrm{new}}$ 具有如下性质：$\phi_E^{\mathrm{new}}(\cdot,J^{\mathrm{new}}\cdot)$ 是对称的；$\pi$ 是 $J^{\mathrm{new}}$-\trHolomorphic{}的；而且在 $TE^v$ 上 $J^{\mathrm{new}}=J$。在这些术语下，\eqref{eq:section-equation} 不过是要求一个\trSection{}在通常意义下为 $J^{\mathrm{new}}$-\trHolomorphic{}的方程。\trConnection{}这一观点在理解解 $u$ 的某些基本性质时也很有用。例如，\trEnergy{}的\trApriori{}估计为
\begin{equation*}
E(u)=\int_S\tfrac12|(Du)^v-Y(u)|^2=\int_{\partial S}\alpha_F+\theta_S-\int_S u^*R,
\end{equation*}
其中 $(\cdot)^v$ 是关于 $TE=TE^h\oplus TE^v$ 的\trVertical{}分量，而\trCurvature{}项在局部\trTrivialization{} \eqref{eq:local-trivialization-forms} 中由
\begin{equation}\label{eq:curvature-term}
R(\partial_s,\partial_t)=\bigl(\partial_s(K-\kappa)(\partial_t)-\partial_t(K-\kappa)(\partial_s)+\{(K-\kappa)(\partial_s),(K-\kappa)(\partial_t)\}\bigr)\,ds\wedge dt.
\end{equation}
给出。\trEquivalence{}地，可以写成 $E(u)=\int_S\tfrac12|(Du)^{v,\mathrm{new}}|^2$，其中使用的是修改后的\trConnection{} \eqref{eq:modified-connection}，而此时 \eqref{eq:curvature-term} 的确就是该\trConnection{}的\trCurvature。
\end{remark}
\begin{remark}\label{thm:moving-boundary-analogue}
在这一语境中也有\trMovingBoundaryCondition{}的一个类似物。这样的\trBoundaryCondition{}无非由任意一族\trExact\trLagrange\trSubmanifold{}$F_z\subset E_z$，$z\in\partial S$ 组成。\trRelative\trPerturbationDatum{}的相应概念仍是一对 $(K,J)$，只是 $K$ 不再在限制到 $F$ 上时为零，而是具有如下性质：对任何 $X\in TF|F_z$，$z\in\partial S$，只要 $D\pi(X)$ 平行于边界，就有
\begin{equation}\label{eq:moving-relative-datum}
d(i_XK)|F_z=\phi(\cdot,X)|F_z.
\end{equation}
除此之外，\trInhomogeneous\trPseudoholomorphic\trSection{}的定义保持不变。其含义应当是清楚的：诸 $F_z$ 未必被 $E$ 上给定\trSymplectic\trConnection{}的\trParallelTransport{}所保持，但换成满足 \eqref{eq:moving-relative-datum} 的某个 $K$ 而得到的修改形式 \eqref{eq:modified-form} 恰好恢复了这种\trCompatible{}性。
\end{remark}
记 $\scr{M}_{E/S}$ 为 \eqref{eq:section-equation} 的全体解所成的\trModuliSpace。该方程在这个\trModuliSpace{}的任一点 $u$ 处的\trLinearization{}由一个 $\bar\partial$-算子 $D_{E/S,u}$ 给出，它定义在\trVectorBundle{}$u^*(TE^v)$ 上，其边界值在 $\partial S$ 上由 $u^*(TF^v)$ 给出。照常，如果 $D_{E/S,u}$ 对一切 $u$ 都是满的，我们就称 $(K,J)$ 是\trRegular{}的，此时该\trModuliSpace{}是一个\trClosed\trSmooth\trManifold。这个性质对\trGeneric{}选取的 $(K,J)$ 成立。取定任一这样的选择，便可通过对其\trIsolated{}点作 mod 2 计数来定义一个玩具模型\trInvariant：
\begin{equation}\label{eq:toy-invariant}
\Phi_{E/S}=\#\scr{M}_{E/S}\in\ZZ/2.
\end{equation}
\begin{example}\label{thm:local-model-sections}
设 $\pi\colon E\to S$ 是\cref{thm:basic-local-model}与\cref{thm:toy-model-thimble}中的局部模型。我们在其上放置一个\trLagrange\trBoundaryCondition{}$F$，它是诸标准\trVanishingCycle{}$F_z=\sqrt{\bar z}\,S^n$（$|z|=r$）的并。在 $F$ 上我们有 $x_j=(\pi(x)/r)\bar x_j$，于是
\begin{equation*}
\theta_E|F=\frac{i}{4}\sum_j x_j\,d\bar x_j-\bar x_j\,dx_j=\pi^*\Bigl(\frac{i}{8r}(z\,d\bar z-\bar z\,dz)\Bigr).
\end{equation*}
定义 $\alpha_F$ 使得 $\alpha_F+\theta_S|\partial S$ 与上式右边一致，并取 $h_F=0$，便表明 $F$ 是一个\trLagrange\trBoundaryCondition（当然，我们本来就知道这会成立，因为诸 $\sqrt{\bar z}\,S^n$ 在任何方向上都保持于\trParallelTransport{}之下）。我们取标准\trComplexStructure{}$J=I_E$ 与平凡的\trInhomogeneous{}项 $K=0$。所得\trModuliSpace{}$\scr{M}_{E/S}$ 可以显式确定：它由映射
\begin{equation}\label{eq:local-model-sections}
u_a\colon S\longrightarrow E\subset\CC^{n+1},\qquad u_a(z)=r^{-1/2}az+r^{1/2}\bar a
\end{equation}
组成，其中 $a\in\CC^{n+1}$ 满足 $a_1^2+\cdots+a_{n+1}^2=0$ 与 $\|a\|^2=1/2$ [131, Lemma 2.16]。此外，这些点都是该\trModuliSpace{}的\trRegular{}点。严格说来，还须稍加小心，即 \eqref{eq:basic-model-domain} 中的 $s$ 应取得足够大，使这些\trSection{}的像的确含于 $E$ 之中。这与 $\partial^hE$ 不是\trHolomorphic\trConvex{}的这一事实有关，正是它使 $E$ 无法成为一个完全合格的\trExact\trLefschetzFibration{}的例子。暂且不计这一缺陷，我们发现所有这些\trSection{}在\trInvariant{} \eqref{eq:toy-invariant} 之下会因维数原因是平凡的（\trModuliSpace{}中没有\trIsolated{}点）。
\end{example}
这一计算为更一般的消没结果提供了原型。设 $\pi\colon E\to S$ 是一个\trLefschetzFibration，其基 $S$ 是一个\trDisk，并且只有一个\trCriticalPoint。在 $\partial S$ 上取一点 $z$，并取一条从 $z$ 通到\trCriticalValue{}的\trVanishingPath{}$\gamma$。我们写 $M=E_z$ 与 $V=V_\gamma$。由 \eqref{eq:symplectic-picard-lefschetz}，绕边界\trLoop{}$\lambda$ 一周的\trMonodromy{}$h_\lambda$ \trIsotopy{}于 $\tau_V$（其方式在固定端点的\trHomotopy{}意义下唯一）。这意味着我们可以按如下方式在我们的\trFibration{}上放置一个\trMovingBoundaryCondition{}$F$：从 $F_z=V$ 出发，用\trParallelTransport{}使它沿边界移动几乎一整圈，然后用\trIsotopy{}$h_\lambda(V)\simeq\tau_V(V)=V$ 把所得结果接回 $F_z$。我们将把这称为一个\emph{标准\trBoundaryCondition}（standard boundary condition）。
\begin{lemma}\label{thm:phi-es-vanishing}
$\Phi_{E/S}=0$.
\end{lemma}
我们只对其证明作一勾画；细节见 [131, Section 2.3]。为简化记号，设 $S=D$ 是单位\trDisk，并且 $\mathrm{Critv}(\pi)=\{0\}$。经过一个局部的\trDeformation（它不影响\trInvariant），可以设在\trCriticalPoint{}周围存在\trHolomorphic\trMorse{}坐标，使\trSymplecticForm{}在其中是标准的。考察由把给定的这个\trLefschetzFibration{}限制到较小的\trDisk{}$S^r=rD$ 上而得到的\trLefschetzFibration{}族 $\pi^r\colon E^r\to S^r$，$r\in(0;1]$。从 $F$ 出发，可以赋予它们一个\trSmooth{}的\trBoundaryCondition{}族 $F^r$，当 $r$ 小时它等于\cref{thm:local-model-sections}的局部模型中所描述的那一个。取 $(E^r,F^r)$ 的一族\trPerturbationDatum，当 $r$ 小时它是平凡的。此时，一个\trCompactness{}论证 [131, Lemma 2.15] 表明：在极限 $r\to0$ 下，$\scr{M}_{E^r/S^r}$ 的所有元素都将落入\trCriticalPoint{}的一个任意小的事先指定的邻域，于是我们便处于与\cref{thm:local-model-sections}相同的情形。最后，一个常规的\trCobordism{}论证把结果带回到原来的 $(E,F)$。
作为补充，我们愿意提及一个更初等但更不一般的证明，它只依赖于诸\trModuliSpace{}的维数。首先我们需要一个关于 $(E,F)$ 的\trTopology{}的基本事实：
\begin{lemma}\label{thm:relative-chern-pullback}
若\trRelative\trFirstChernClass{}$2c_1(M,V)\in H^2(M,V)$ 为零，则其对应物 $2c_1(E,F)\in H^2(E,F)$ 是 $H^2(S,\partial S)$ 的标准\trGenerator{}的 $(n+1)$ 倍的\trPullback。
\end{lemma}
由\cref{sub:relative-volume-forms}我们知道，$E$ 容许一个二次复\trVolume{}形式，从而 $2c_1(E)=0$，这意味着 $2c_1(E,F)$ 来自 $H^1(F)$ 的某个元素。当 $n>1$ 时，这意味着它是 $H^2(S,\partial S)$ 中某个元素在 $\pi^*$ 下的像。当 $n=1$ 时，注意到 $2c_1(E,F)$ 在 $H^2(\Delta,V)$ 上的限制必为零——这是由于 $\Delta$（即\trLefschetzThimble）的\trLagrange{}性质——便可得出同样的结论。记住了这一点，论证便归结为一次可以在局部模型中完成的计算，我们将其省略。
\trModuliSpace{}$\scr{M}_{E/S}$ 在点 $u$ 处的期望维数由算子 $D_{E/S,u}$ 的\trIndex{}给出。应用相对\trRiemann--\trRoch{}公式（\cref{thm:det-line-bundle-loop-space}，表述稍有不同）与\cref{thm:relative-chern-pullback}，便发现
\begin{equation*}
\mathrm{index}(D_{E/S,u})=n+\int_S2c_1(u^*(TE^v),u^*(TF^v))=2n-1,
\end{equation*}
这当然与我们在局部模型中所见的相符。特别地，无论怎样选取（\trRegular{}的）\trRelative\trPerturbationDatum，$\scr{M}_{E/S}$ 中都绝不会有\trIsolated{}点。这与前一个（更一般的）\trCobordism{}论证的结论形成对照，后者只保证了这类\trIsolated{}点之间的相互抵消。
